% Part: first-order-logic
% Chapter: beyond
% Section: higher-order-logic

\documentclass[../../../include/open-logic-section]{subfiles}

\begin{document}

\olfileid{fol}{byd}{hol}

\olsection{Logika Orde Luwih Dhuwur}

Saka logika orde kapisan menyang logika orde kapindho,
kita bisa ngrembug himpunan objek ing !!{domain}
orde kapisan kanthi basa formal. Apa sebabé kudu
mandheg ing kono? Upamane, logika orde katelu
ngidini kita ngrembug himpunan saka himpunan objek,
utawa malah himpunan kang ngemot objek lan himpunan
objek bebarengan. Logika orde kapapat banjur
ngidini kita ngrembug himpunan objek kang kaya
mangkono. Gagasan iki bisa diterusake nganti
sembarang tingkat.

Ing praktik, !!{formula} logika orde luwih dhuwur
asring katulis nganggo fungsi tinimbang relasi.
(Yen panyamaan alamiahe digatekake, bedane iki
ora wigati.) Yen ana sawetara "jinis" dhasar
$A$, $B$, $C$,~\dots (saiki kita arani "tipe"),
kita bisa nggawe tipe anyar kanthi aturan
\begin{quote}
Yen $\sigma$ lan $\tau$ iku tipe winates, mula
$\sigma \to \tau$ uga tipe winates.
\end{quote}
Tipe bisa dianggep "label" sintaktis kang
nggolongake objek kang dikarepake ing !!{domain};
$\sigma \to \tau$ njlentrehake objek awujud
fungsi kang nggawa objek tipe~$\sigma$ menyang
objek tipe~$\tau$. Upamane, kita bisa nduweni
tipe~$\Omega$ kanggo nilai bebener, "bener" lan
"kleru", sarta tipe~$\Nat$ kanggo wilangan asli.
Ing kene, objek tipe $\Nat \to \Omega$ bisa
dianggep relasi unari utawa subhimpunan
saka~$\Nat$; objek tipe $\Nat \to \Nat$ iku
fungsi saka wilangan asli menyang wilangan asli;
lan objek tipe $(\Nat \to \Nat) \to \Nat$
iku "fungsional", yaiku fungsi kanthi tipe
luwih dhuwur kang nggawa fungsi menyang wilangan.

Kaya ing logika orde kapindho, logika orde luwih
dhuwur bisa dianggep sawijining logika akeh
jinis, kanthi siji jinis kanggo saben tipe
objek kang arep ditliti. Nanging lumrahe
luwih cetha yen sintaks logika tipe luwih
dhuwur ditemtokake saka dhasare. Upamane,
himpunan tipe winates bisa ditemtokake
kanthi induksi kaya mangkene:
\begin{enumerate}
\item $\Nat$ iku tipe winates.
\item Yen $\sigma$ lan $\tau$ iku tipe winates,
  mula $\sigma \to \tau$ uga tipe winates.
\item Yen $\sigma$ lan $\tau$ iku tipe winates,
  mula $\sigma \times \tau$ uga tipe winates.
\end{enumerate}
Miturut pangertèn lumrah, $\Nat$ nuduhake
tipe wilangan asli, $\sigma \to \tau$ nuduhake
tipe fungsi saka $\sigma$ menyang $\tau$,
lan $\sigma \times \tau$ nuduhake tipe pasangan
objek, siji saka $\sigma$ lan siji saka $\tau$.
Sabanjure, himpunan term bisa ditemtokake
kanthi induksi kaya mangkene:
\begin{enumerate}
\item Kanggo saben tipe $\sigma$, ana variabel
  $x$, $y$, $z$, \dots kanthi tipe $\sigma$
\item $\Obj 0$ iku term tipe $\Nat$
\item $\Obj S$ (penerus) iku term tipe
  $\Nat \to \Nat$
\item Yen $s$ iku term tipe $\sigma$, lan
  $t$ iku term tipe $\Nat \to (\sigma \to \sigma)$,
  mula $\Obj{R}_{st}$ iku term tipe
  $\Nat \to \sigma$
\item Yen $s$ iku term tipe $\tau \to \sigma$
  lan $t$ iku term tipe~$\tau$, mula $s(t)$
  iku term tipe $\sigma$
\item Yen $s$ iku term tipe~$\sigma$ lan $x$
  iku variabel tipe~$\tau$, mula
  $\lambd[x][s]$ iku term tipe
  $\tau \to \sigma$.
\item Yen $s$ iku term tipe~$\sigma$ lan $t$
  iku term tipe~$\tau$, mula $\tuple{s, t}$
  iku term tipe $\sigma \times \tau$.
\item Yen $s$ iku term tipe~$\sigma \times \tau$,
  mula $p_1(s)$ iku term tipe~$\sigma$ lan
  $p_2(s)$ iku term tipe~$\tau$.
\end{enumerate}
Miturut pangertèn lumrah, $\Obj{R}_{st}$
nuduhake fungsi kang ditemtokake kanthi
rekursi dening
\begin{align*}
\Obj{R}_{st}(0) & = s \\
\Obj{R}_{st}(x+1) & = t(x, R_{st}(x)),
\end{align*}
$\tuple{s, t}$ nuduhake pasangan kang
komponen kapisane~$s$ lan komponen kapindhone~$t$,
dene $p_1(s)$ lan~$p_2(s)$ nuduhake unsur
kapisan lan kapindho ("proyeksi") saka~$s$.
Pungkasane, $\lambd[x][s]$ nuduhake
fungsi~$f$ kang ditemtokake dening
\[
f(x) = s
\]
kanggo saben~$x$ kanthi tipe~$\tau$;
mula butir (6) menehi salah siji wujud
komprehensi, saéngga kita bisa nemtokake
fungsi nganggo term.
Cathetan koreksi sumber OLPL-085: ing butir (6)
variabel x nduweni tipe tau, dudu sigma kaya
kang katulis ing prosa sumber sakwise rumus iki.
!!^{formula}s dibangun saka ukara !!{identity}
$\eq[s][t]$ antarane term kanthi tipe padha,
panyambung proposisional lumrah, lan
kuantifikasi tumrap tipe luwih dhuwur.
Aksioma sistem banjur bisa dijupuk minangka
persamaan dhasar kang ngatur term-term ing
ndhuwur, bebarengan karo aturan logika
lumrah kanggo kuantor lan !!{identity}.

Yen sistem tipe winates ditambahi tipe~$\Omega$
kanggo nilai bebener, kita uga kudu nambah
aksioma kang ngatur panganggone. Malah yen
kita cukup trampil, !!{formula}s kang rumit
bisa dibusak kabeh lan diganti term tipe~$\Omega$!{}
Sistem bukti banjur bisa diowahi
manut kahanan iku. Asile yaiku, ing pokoke,
\emph{teori tipe prasaja} kang diajokake
Alonzo Church ing dasawarsa 1930-an.

Kaya ing logika orde kapindho, semantik
tipe luwih dhuwur nduweni sawetara versi
kang bisa dipilih. Ing versi pepak,
variabel tipe $\sigma \to \tau$ nglimputi
himpunan \emph{kabeh} fungsi saka objek
tipe~$\sigma$ menyang objek tipe~$\tau$.
Kaya kang bisa dikira, semantik iki
kuwat banget nganti ora ana sistem
!!{derivation} efektif kang jangkep
kanggo semantik kasebut. Nanging kita
bisa nggatekake semantik kang luwih ringkih,
ing ngendi !!a{structure} dumadi saka
himpunan unsur $T_\tau$ kanggo saben tipe
$\tau$, bebarengan karo operasi kang
cocog kanggo aplikasi, proyeksi, lan
sapanunggalane. Yen rinciane ditata
kanthi bener, teorema kelengkapan
bisa dioleh kanggo jinis sistem
!!{derivation} kang diterangake ing ndhuwur.

Logika tipe luwih dhuwur narik kawigaten
amarga nyedhiyakake kerangka kang bisa
ngemot akeh matematika kanthi cara lumrah:
wiwit saka~$\Nat$, kita bisa nemtokake
wilangan real, fungsi kontinu, lan liyane.
Logika iki uga narik kawigaten mligine
ing logika intuisionistik, amarga tipe-tipe
kasebut nduweni tafsiran "konstruktif"
kang cetha. Malah versi konstruktif
semantik tipe luwih dhuwur bisa
dikembangake adhedhasar logika intuisionistik,
dudu logika klasik; versi iki njlentrehake
tafsiran konstruktif kasebut kanthi
cetha lan ing akeh bab luwih narik
kawigaten tinimbang pasangan klasiké.

\end{document}
